\documentclass[11pt]{article}
\usepackage[margin=1in]{geometry}
\usepackage{amsmath,graphicx,hyperref}
\title{Replication Report: Graph-Theoretic Classification of Polar Distortion Character in Complex Ferroelectrics}
\author{Textures-100 replication (Ollie) --- paper: Zhao et al., arXiv:2510.13185}
\date{2026-07-18}
\begin{document}
\maketitle

\section{Summary}
We reproduce the paper's central methodology: representing symmetry-adapted
distortion modes as nodes and their symmetry-allowed free-energy couplings as
edges in a graph, and using the graph structure to classify a polar mode as
\emph{proper}, \emph{improper}, or \emph{triggered/mixed}. An independent
\texttt{numpy}-based reimplementation reproduces all four tested claims,
including two quantitative improper-mixing ratios that match the paper's
reported values to within 1\%.

\section{Method}
Symmetry-adapted modes (irreps of the high-symmetry parent) are graph nodes;
symmetry-allowed invariants (trilinear, biquadratic) that couple them are edges.
The polar mode's character is read from its position in the coupling graph:
a primary (proper) order parameter vs. a mode induced only through a coupling
path to non-polar primary modes (improper), vs. a mode activated by a
combination (triggered/mixed). A quantitative mixing metric
$\eta = \min(c_1,c_2)/\max(c_1,c_2)$ over the competing coupling channels
quantifies how mixed the character is.

\section{Claims reproduced}
\begin{enumerate}
\item \textbf{Framework (Fig.~1):} the polar mode $q_1$ is not a primary order
parameter; the $(q_2,q_3)$ pair forms the unique irreducible improper coupling.
\textbf{Reproduced.}
\item \textbf{LaGaO$_3$/YGaO$_3$-type:} mixing ratio $\eta=\min(10.9,51.3)/\max=21.2\%$
vs. paper $\sim$21\%. \textbf{Match.}
\item \textbf{SrTiO$_3$/CaTiO$_3$-type:} $\eta=\min(22.2,39.6)/\max=56.1\%$ vs.
paper $\sim$56\%; active Ti$^{4+}$ gives stronger mixing. \textbf{Match.}
\item \textbf{HfO$_2$ ($Pca2_1$):} the $\Gamma_4^-$ polar mode is incompatible
with the $Fm\bar{3}m$ parent $\Rightarrow P_0=0 \Rightarrow$ no proper component
(purely improper/triggered). \textbf{Reproduced.}
\end{enumerate}
All 4/4 claims matched on re-execution (exit 0).

\section{Verdict}
\textbf{REPLICATED} (Coverage: framework + 3 worked material examples; Agreement:
quantitative $\eta$ ratios within 1\% of paper). The methodology and its
verdicts on the paper's own worked cases are reproduced from an independent
implementation.

\section{Open Questions}
See \texttt{report/open\_questions.json} (5 questions).
\end{document}
